% Generated by scripts/11_paper.py from paper/references.json (records fetched from Crossref or arXiv
% by scripts/11c_references.py). Do not edit: change paper/references.yaml and run both scripts.
{\small\begin{thebibliography}{81}\setlength{\itemsep}{1pt}
\bibitem{kang2018peerread} Kang D, Ammar W, Dalvi B, Zuylen M, Kohlmeier S, Hovy E, et al. A Dataset of Peer Reviews (PeerRead): Collection, Insights and NLP Applications. In: Proceedings of the 2018 Conference of the North American Chapter of the Association for Computational Linguistics: Human Language Technologies, Volume 1 (Long Papers); 2018. p. 1647-1661. doi: \href{https://doi.org/10.18653/v1/n18-1149}{\nolinkurl{10.18653/v1/n18-1149}}
\bibitem{wang2023openreview} Wang G, Peng Q, Zhang Y, Zhang M. What have we learned from OpenReview? World Wide Web. 2023;26(2):683-708. doi: \href{https://doi.org/10.1007/s11280-022-01109-z}{\nolinkurl{10.1007/s11280-022-01109-z}}
\bibitem{shah2022challenges} Shah NB. Challenges, experiments, and computational solutions in peer review. Communications of the ACM. 2022;65(6):76-87. doi: \href{https://doi.org/10.1145/3528086}{\nolinkurl{10.1145/3528086}}
\bibitem{ghosal2019deepsentipeer} Ghosal T, Verma R, Ekbal A, Bhattacharyya P. DeepSentiPeer: Harnessing Sentiment in Review Texts to Recommend Peer Review Decisions. In: Proceedings of the 57th Annual Meeting of the Association for Computational Linguistics; 2019. p. 1120-1130. doi: \href{https://doi.org/10.18653/v1/p19-1106}{\nolinkurl{10.18653/v1/p19-1106}}
\bibitem{deng2020habnet} Deng Z, Peng H, Xia C, Li J, He L, Yu P. Hierarchical Bi-Directional Self-Attention Networks for Paper Review Rating Recommendation. In: Proceedings of the 28th International Conference on Computational Linguistics; 2020. p. 6302-6314. doi: \href{https://doi.org/10.18653/v1/2020.coling-main.555}{\nolinkurl{10.18653/v1/2020.coling-main.555}}
\bibitem{bharti2021peerassist} Bharti PK, Ranjan S, Ghosal T, Agrawal M, Ekbal A. PEERAssist: Leveraging on Paper-Review Interactions to Predict Peer Review Decisions. In: Towards Open and Trustworthy Digital Societies. Lecture Notes in Computer Science; 2021. p. 421-435. doi: \href{https://doi.org/10.1007/978-3-030-91669-5_33}{\nolinkurl{10.1007/978-3-030-91669-5_33}}
\bibitem{wang2018sentiment} Wang K, Wan X. Sentiment Analysis of Peer Review Texts for Scholarly Papers. In: The 41st International ACM SIGIR Conference on Research \& Development in Information Retrieval; 2018. p. 175-184. doi: \href{https://doi.org/10.1145/3209978.3210056}{\nolinkurl{10.1145/3209978.3210056}}
\bibitem{schlichtkrull2018rgcn} Schlichtkrull M, Kipf TN, Bloem P, van den Berg R, Titov I, Welling M. Modeling Relational Data with Graph Convolutional Networks. In: The Semantic Web. Lecture Notes in Computer Science; 2018. p. 593-607. doi: \href{https://doi.org/10.1007/978-3-319-93417-4_38}{\nolinkurl{10.1007/978-3-319-93417-4_38}}
\bibitem{wang2019han} Wang X, Ji H, Shi C, Wang B, Ye Y, Cui P, et al. Heterogeneous Graph Attention Network. In: The World Wide Web Conference; 2019. p. 2022-2032. doi: \href{https://doi.org/10.1145/3308558.3313562}{\nolinkurl{10.1145/3308558.3313562}}
\bibitem{hu2020hgt} Hu Z, Dong Y, Wang K, Sun Y. Heterogeneous Graph Transformer. In: Proceedings of The Web Conference 2020; 2020. p. 2704-2710. doi: \href{https://doi.org/10.1145/3366423.3380027}{\nolinkurl{10.1145/3366423.3380027}}
\bibitem{ying2019gnnexplainer} Ying R, Bourgeois D, You J, Zitnik M, Leskovec J. GNNExplainer: Generating Explanations for Graph Neural Networks. arXiv:1903.03894 [Preprint]. 2019. doi: \href{https://doi.org/10.48550/arxiv.1903.03894}{\nolinkurl{10.48550/arxiv.1903.03894}}
\bibitem{yuan2023explainability} Yuan H, Yu H, Gui S, Ji S. Explainability in Graph Neural Networks: A Taxonomic Survey. IEEE Transactions on Pattern Analysis and Machine Intelligence. 2022:1-19. doi: \href{https://doi.org/10.1109/tpami.2022.3204236}{\nolinkurl{10.1109/tpami.2022.3204236}}
\bibitem{kapoor2023leakage} Kapoor S, Narayanan A. Leakage and the reproducibility crisis in machine-learning-based science. Patterns. 2023;4(9):100804. doi: \href{https://doi.org/10.1016/j.patter.2023.100804}{\nolinkurl{10.1016/j.patter.2023.100804}}
\bibitem{wen2026decoupling} Wen Y. Decoupling Scores and Text: The Politeness Principle in Peer Review. arXiv:2604.14162 [Preprint]. 2026. doi: \href{https://doi.org/10.48550/arxiv.2604.14162}{\nolinkurl{10.48550/arxiv.2604.14162}}
\bibitem{li2026reviewgraph} Li S, Fan L, Lin Y, Li Z, Wei X, Ni S, et al. Automatic Paper Reviewing with Heterogeneous Graph Reasoning over LLM-Simulated Reviewer-Author Debates. Proceedings of the AAAI Conference on Artificial Intelligence. 2026;40(37):31717-31725. doi: \href{https://doi.org/10.1609/aaai.v40i37.40439}{\nolinkurl{10.1609/aaai.v40i37.40439}}
\bibitem{lin2023moprd} Lin J, Song J, Zhou Z, Chen Y, Shi X. MOPRD: A multidisciplinary open peer review dataset. Neural Computing and Applications. 2023;35(34):24191-24206. doi: \href{https://doi.org/10.1007/s00521-023-08891-5}{\nolinkurl{10.1007/s00521-023-08891-5}}
\bibitem{bao2021decisionsets} Bao P, Hong W, Li X. Predicting Paper Acceptance via Interpretable Decision Sets. In: Companion Proceedings of the Web Conference 2021; 2021. p. 461-467. doi: \href{https://doi.org/10.1145/3442442.3451370}{\nolinkurl{10.1145/3442442.3451370}}
\bibitem{rimkus2024reviewer} Rimkus L, Verbickas J, Batista-Navarro R. Reviewer 2 Must Be Stopped: Transformer-Based Approaches for Predicting Paper Acceptance. In: Information Management and Big Data. Communications in Computer and Information Science; 2024. p. 249-260. doi: \href{https://doi.org/10.1007/978-3-031-63616-5_19}{\nolinkurl{10.1007/978-3-031-63616-5_19}}
\bibitem{pradhan2021mrgen} Pradhan T, Bhatia C, Kumar P, Pal S. A deep neural architecture based meta-review generation and final decision prediction of a scholarly article. Neurocomputing. 2021;428:218-238. doi: \href{https://doi.org/10.1016/j.neucom.2020.11.004}{\nolinkurl{10.1016/j.neucom.2020.11.004}}
\bibitem{ribeiro2021acceptance} Ribeiro AC, Sizo A, Lopes Cardoso H, Reis LP. Acceptance Decision Prediction in Peer-Review Through Sentiment Analysis. In: Progress in Artificial Intelligence. Lecture Notes in Computer Science; 2021. p. 766-777. doi: \href{https://doi.org/10.1007/978-3-030-86230-5_60}{\nolinkurl{10.1007/978-3-030-86230-5_60}}
\bibitem{prama2024classification} Prama TT, Biswas AA, Anwar MM. Deep Learning-Based Classification of Conference Paper Reviews: Accept or Reject? In: Intelligent Systems Design and Applications. Lecture Notes in Networks and Systems; 2024. p. 64-74. doi: \href{https://doi.org/10.1007/978-3-031-64776-5_7}{\nolinkurl{10.1007/978-3-031-64776-5_7}}
\bibitem{kumar2026reviews} Kumar S, Arora H, Ghosal T, Ekbal A. From reviews to decisions: A joint multitask aspect sentiment leveraged framework for assisting decision prediction from academic peer reviews. Scientometrics. 2026;131(1):583-608. doi: \href{https://doi.org/10.1007/s11192-025-05529-2}{\nolinkurl{10.1007/s11192-025-05529-2}}
\bibitem{muangkammuen2022semi} Muangkammuen P, Fukumoto F, Li J, Suzuki Y. Exploiting Labeled and Unlabeled Data via Transformer Fine-tuning for Peer-Review Score Prediction. In: Findings of the Association for Computational Linguistics: EMNLP 2022; 2022. p. 2233-2240. doi: \href{https://doi.org/10.18653/v1/2022.findings-emnlp.164}{\nolinkurl{10.18653/v1/2022.findings-emnlp.164}}
\bibitem{muangkammuen2023transfer} Muangkammuen P, Fukumoto F, Li J, Suzuki Y. Intermediate-Task Transfer Learning for Peer Review Score Prediction. In: Proceedings of the 13th International Joint Conference on Natural Language Processing and the 3rd Conference of the Asia-Pacific Chapter of the Association for Computational Linguistics: Student Research Workshop; 2023. p. 40-47. doi: \href{https://doi.org/10.18653/v1/2023.ijcnlp-srw.6}{\nolinkurl{10.18653/v1/2023.ijcnlp-srw.6}}
\bibitem{fytas2021accepted} Fytas P, Rizos G, Specia L. What Makes a Scientific Paper be Accepted for Publication? In: Proceedings of the First Workshop on Causal Inference and NLP; 2021. p. 44-60. doi: \href{https://doi.org/10.18653/v1/2021.cinlp-1.4}{\nolinkurl{10.18653/v1/2021.cinlp-1.4}}
\bibitem{chakraborty2020aspect} Chakraborty S, Goyal P, Mukherjee A. Aspect-based Sentiment Analysis of Scientific Reviews. In: Proceedings of the ACM/IEEE Joint Conference on Digital Libraries in 2020; 2020. p. 207-216. doi: \href{https://doi.org/10.1145/3383583.3398541}{\nolinkurl{10.1145/3383583.3398541}}
\bibitem{konan2026paperscorer} Konan S, Liu J, Liu Z. PaperScorer: How Predictable Is Paper Acceptance? arXiv:2610.31791 [Preprint]. 2026. doi: \href{https://doi.org/10.48550/arxiv.2610.31791}{\nolinkurl{10.48550/arxiv.2610.31791}}
\bibitem{heo2026drives} Jung S, Pyeon G, Heo I, Ahn H. What Drives Paper Acceptance? A Process-Centric Analysis of Modern Peer Review. arXiv:2509.25701 [Preprint]. 2025. doi: \href{https://doi.org/10.48550/arxiv.2509.25701}{\nolinkurl{10.48550/arxiv.2509.25701}}
\bibitem{hua2019argument} Hua X, Nikolov M, Badugu N, Wang L. Argument Mining for Understanding Peer Reviews. In: Proceedings of the 2019 Conference of the North American Chapter of the Association for Computational Linguistics: Human Language Technologies, Volume 1 (Long and Short Papers); 2019. p. 2131-2137. doi: \href{https://doi.org/10.18653/v1/n19-1219}{\nolinkurl{10.18653/v1/n19-1219}}
\bibitem{fromm2021argument} Fromm M, Faerman E, Berrendorf M, Bhargava S, Qi R, Zhang Y, et al. Argument Mining Driven Analysis of Peer-Reviews. Proceedings of the AAAI Conference on Artificial Intelligence. 2021;35(6):4758-4766. doi: \href{https://doi.org/10.1609/aaai.v35i6.16607}{\nolinkurl{10.1609/aaai.v35i6.16607}}
\bibitem{guo2023substantiation} Guo Y, Shang G, Rennard V, Vazirgiannis M, Clavel C. Automatic Analysis of Substantiation in Scientific Peer Reviews. In: Findings of the Association for Computational Linguistics: EMNLP 2023; 2023. p. 10198-10216. doi: \href{https://doi.org/10.18653/v1/2023.findings-emnlp.684}{\nolinkurl{10.18653/v1/2023.findings-emnlp.684}}
\bibitem{yuan2022automate} Yuan W, Liu P, Neubig G. Can We Automate Scientific Reviewing? Journal of Artificial Intelligence Research. 2022;75:171-212. doi: \href{https://doi.org/10.1613/jair.1.12862}{\nolinkurl{10.1613/jair.1.12862}}
\bibitem{bharti2025multipeer} Bharti PK, Agarwal M, Ekbal A. MultiPEER: A deep, multi-task framework for identifying aspect categories and their sentiments in scientific peer reviews. Journal of Information Science. 2025:01655515251353189. doi: \href{https://doi.org/10.1177/01655515251353189}{\nolinkurl{10.1177/01655515251353189}}
\bibitem{kuznetsov2022revise} Kuznetsov I, Buchmann J, Eichler M, Gurevych I. Revise and Resubmit: An Intertextual Model of Text-based Collaboration in Peer Review. Computational Linguistics. 2022;48(4):949-986. doi: \href{https://doi.org/10.1162/coli_a_00455}{\nolinkurl{10.1162/coli_a_00455}}
\bibitem{thelwall2020praise} Thelwall M, Papas ER, Nyakoojo Z, Allen L, Weigert V. Automatically detecting open academic review praise and criticism. Online Information Review. 2020;44(5):1057-1076. doi: \href{https://doi.org/10.1108/oir-11-2019-0347}{\nolinkurl{10.1108/oir-11-2019-0347}}
\bibitem{li2022capsule} Li S, Li Y, Zhao WX, Ding B, Wen JR. Interpretable Aspect-Aware Capsule Network for Peer Review Based Citation Count Prediction. ACM Transactions on Information Systems. 2022;40(1):1-29. doi: \href{https://doi.org/10.1145/3466640}{\nolinkurl{10.1145/3466640}}
\bibitem{sun2024textual} Sun Z. Textual features of peer review predict top-cited papers: An interpretable machine learning perspective. Journal of Informetrics. 2024;18(2):101501. doi: \href{https://doi.org/10.1016/j.joi.2024.101501}{\nolinkurl{10.1016/j.joi.2024.101501}}
\bibitem{gao2019rebuttal} Gao Y, Eger S, Kuznetsov I, Gurevych I, Miyao Y. Does My Rebuttal Matter? Insights from a Major NLP Conference. In: Proceedings of the 2019 Conference of the North American Chapter of the Association for Computational Linguistics: Human Language Technologies, Volume 1 (Long and Short Papers); 2019. p. 1274-1290. doi: \href{https://doi.org/10.18653/v1/n19-1129}{\nolinkurl{10.18653/v1/n19-1129}}
\bibitem{cheng2020ape} Cheng L, Bing L, Yu Q, Lu W, Si L. APE: Argument Pair Extraction from Peer Review and Rebuttal via Multi-task Learning. In: Proceedings of the 2020 Conference on Empirical Methods in Natural Language Processing (EMNLP); 2020. p. 7000-7011. doi: \href{https://doi.org/10.18653/v1/2020.emnlp-main.569}{\nolinkurl{10.18653/v1/2020.emnlp-main.569}}
\bibitem{kennard2022disapere} Kennard N, O’Gorman T, Das R, Sharma A, Bagchi C, Clinton M, et al. DISAPERE: A Dataset for Discourse Structure in Peer Review Discussions. In: Proceedings of the 2022 Conference of the North American Chapter of the Association for Computational Linguistics: Human Language Technologies; 2022. p. 1234-1249. doi: \href{https://doi.org/10.18653/v1/2022.naacl-main.89}{\nolinkurl{10.18653/v1/2022.naacl-main.89}}
\bibitem{huang2023rebuttal} Huang J, Huang WB, Bu Y, Cao Q, Shen H, Cheng X. What makes a successful rebuttal in computer science conferences?: A perspective on social interaction. Journal of Informetrics. 2023;17(3):101427. doi: \href{https://doi.org/10.1016/j.joi.2023.101427}{\nolinkurl{10.1016/j.joi.2023.101427}}
\bibitem{cai2026score} Cai S, Wang X. Score increases, author–reviewer communication, and subsequent citations in open peer review: Evidence from ICLR. Journal of Informetrics. 2026;20(4):101878. doi: \href{https://doi.org/10.1016/j.joi.2026.101878}{\nolinkurl{10.1016/j.joi.2026.101878}}
\bibitem{kargaran2025insights} Kargaran AH, Nikeghbal N, Yang J, Ousidhoum N. Insights from the ICLR Peer Review and Rebuttal Process. arXiv:2511.15462 [Preprint]. 2025. doi: \href{https://doi.org/10.48550/arxiv.2511.15462}{\nolinkurl{10.48550/arxiv.2511.15462}}
\bibitem{zhuang2025survey} Zhuang Z, Chen J, Xu H, Jiang Y, Lin J. Large language models for automated scholarly paper review: A survey. Information Fusion. 2025;124:103332. doi: \href{https://doi.org/10.1016/j.inffus.2025.103332}{\nolinkurl{10.1016/j.inffus.2025.103332}}
\bibitem{zhou2024reliable} Zhou R, Chen L, Yu K. Is LLM a Reliable Reviewer? A Comprehensive Evaluation of LLM on Automatic Paper Reviewing Tasks. In: Proceedings of the 2024 Joint International Conference on Computational Linguistics, Language Resources and Evaluation (LREC-COLING 2024). Proceedings of the Language Resources and Evaluation Conference; 2024. p. 9340-9351. doi: \href{https://doi.org/10.63317/48d359hjdvog}{\nolinkurl{10.63317/48d359hjdvog}}
\bibitem{thelwall2025chatgpt} Thelwall M, Yaghi A. Evaluating the predictive capacity of ChatGPT for academic peer review outcomes across multiple platforms. Scientometrics. 2025;130(10):5285-5307. doi: \href{https://doi.org/10.1007/s11192-025-05287-1}{\nolinkurl{10.1007/s11192-025-05287-1}}
\bibitem{yang2026alignment} Yang Z, Zhou X, Wang H, Jiang Y, Zhang X. Large Language Models in Peer Review: Decision Alignment, Review-Text Characteristics, and Human–AI Aggregation at ICLR 2025. Publications. 2026;14(3):55. doi: \href{https://doi.org/10.3390/publications14030055}{\nolinkurl{10.3390/publications14030055}}
\bibitem{jin2024agentreview} Jin Y, Zhao Q, Wang Y, Chen H, Zhu K, Xiao Y, et al. AgentReview: Exploring Peer Review Dynamics with LLM Agents. In: Proceedings of the 2024 Conference on Empirical Methods in Natural Language Processing; 2024. p. 1208-1226. doi: \href{https://doi.org/10.18653/v1/2024.emnlp-main.70}{\nolinkurl{10.18653/v1/2024.emnlp-main.70}}
\bibitem{liang2024monitoring} Liang W, Izzo Z, Zhang Y, Lepp H, Cao H, Zhao X, et al. Monitoring AI-Modified Content at Scale: A Case Study on the Impact of ChatGPT on AI Conference Peer Reviews. arXiv:2403.07183 [Preprint]. 2024. doi: \href{https://doi.org/10.48550/arxiv.2403.07183}{\nolinkurl{10.48550/arxiv.2403.07183}}
\bibitem{russo2025lottery} Russo G, Horta Ribeiro M, Davidson TR, Veselovsky V, West R. The AI Review Lottery: Widespread AI-Assisted Peer Reviews Boost Paper Scores and Acceptance Rates. Proceedings of the ACM on Human-Computer Interaction. 2025;9(7):1-28. doi: \href{https://doi.org/10.1145/3757667}{\nolinkurl{10.1145/3757667}}
\bibitem{thakkar2026feedback} Thakkar N, Yuksekgonul M, Silberg J, Garg A, Peng N, Sha F, et al. A large-scale randomized study of large language model feedback in peer review. Nature Machine Intelligence. 2026;8(3):326-336. doi: \href{https://doi.org/10.1038/s42256-026-01188-x}{\nolinkurl{10.1038/s42256-026-01188-x}}
\bibitem{ebrahimi2025rottenreviews} Ebrahimi S, Sadeghian S, Ghorbanpour A, Arabzadeh N, Salamat S, Li M, et al. RottenReviews: Benchmarking Review Quality with Human and LLM-Based Judgments. In: Proceedings of the 34th ACM International Conference on Information and Knowledge Management; 2025. p. 5642-5649. doi: \href{https://doi.org/10.1145/3746252.3761506}{\nolinkurl{10.1145/3746252.3761506}}
\bibitem{tang2026quality} Tang Y, Kang Y, Wu S, Zhang R, Sun Z. Can large language models assess the quality of peer review? An empirical study. Scientometrics. 2026;131(4):2237-2259. doi: \href{https://doi.org/10.1007/s11192-026-05622-0}{\nolinkurl{10.1007/s11192-026-05622-0}}
\bibitem{he2026pragmatic} He G, Su T, Zhang M, Yuan J, Jiang Z, Ma L, et al. Toward better pragmatic tagging of peer review: Enhancing benchmark datasets via human-in-the-loop multi-agent collaboration. Information Processing \& Management. 2026;63(6):104704. doi: \href{https://doi.org/10.1016/j.ipm.2026.104704}{\nolinkurl{10.1016/j.ipm.2026.104704}}
\bibitem{kipf2017gcn} Kipf TN, Welling M. Semi-Supervised Classification with Graph Convolutional Networks. arXiv:1609.02907 [Preprint]. 2016. doi: \href{https://doi.org/10.48550/arxiv.1609.02907}{\nolinkurl{10.48550/arxiv.1609.02907}}
\bibitem{cohan2020specter} Cohan A, Feldman S, Beltagy I, Downey D, Weld D. SPECTER: Document-level Representation Learning using Citation-informed Transformers. In: Proceedings of the 58th Annual Meeting of the Association for Computational Linguistics; 2020. p. 2270-2282. doi: \href{https://doi.org/10.18653/v1/2020.acl-main.207}{\nolinkurl{10.18653/v1/2020.acl-main.207}}
\bibitem{singh2023scirepeval} Singh A, D’Arcy M, Cohan A, Downey D, Feldman S. SciRepEval: A Multi-Format Benchmark for Scientific Document Representations. In: Proceedings of the 2023 Conference on Empirical Methods in Natural Language Processing; 2023. p. 5548-5566. doi: \href{https://doi.org/10.18653/v1/2023.emnlp-main.338}{\nolinkurl{10.18653/v1/2023.emnlp-main.338}}
\bibitem{bharti2026consistentpeer} Bharti PK, Panchal M, Dalal V. ConsistentPeer: reviewers through GraphRAG-driven counterfactuals to measure consistency in peer review. International Journal of Data Science and Analytics. 2026;22(1):64. doi: \href{https://doi.org/10.1007/s41060-026-01047-3}{\nolinkurl{10.1007/s41060-026-01047-3}}
\bibitem{xue2022dgcbert} Xue Z, He G, Liu J, Jiang Z, Zhao S, Lu W. Re-examining lexical and semantic attention: Dual-view graph convolutions enhanced BERT for academic paper rating. Information Processing \& Management. 2023;60(2):103216. doi: \href{https://doi.org/10.1016/j.ipm.2022.103216}{\nolinkurl{10.1016/j.ipm.2022.103216}}
\bibitem{sundararajan2017ig} Sundararajan M, Taly A, Yan Q. Axiomatic Attribution for Deep Networks. arXiv:1703.01365 [Preprint]. 2017. doi: \href{https://doi.org/10.48550/arxiv.1703.01365}{\nolinkurl{10.48550/arxiv.1703.01365}}
\bibitem{agarwal2023evaluating} Agarwal C, Queen O, Lakkaraju H, Zitnik M. Evaluating explainability for graph neural networks. Scientific Data. 2023;10(1):144. doi: \href{https://doi.org/10.1038/s41597-023-01974-x}{\nolinkurl{10.1038/s41597-023-01974-x}}
\bibitem{li2023xpath} Li T, Deng J, Shen Y, Qiu L, Yongxiang H, Cao CC. Towards Fine-Grained Explainability for Heterogeneous Graph Neural Network. Proceedings of the AAAI Conference on Artificial Intelligence. 2023;37(7):8640-8647. doi: \href{https://doi.org/10.1609/aaai.v37i7.26040}{\nolinkurl{10.1609/aaai.v37i7.26040}}
\bibitem{jiang2026metapath} Jiang M, Ali N, Ding Y, Zhang Y. Is Meta-Path Attention an Explanation? Evidence of Alignment and Decoupling in Heterogeneous GNNs. arXiv:2602.08500 [Preprint]. 2026. doi: \href{https://doi.org/10.48550/arxiv.2602.08500}{\nolinkurl{10.48550/arxiv.2602.08500}}
\bibitem{cvetkovic2026explainability} Cvetković D, Mitrović Dankulov M. Assessing the explainability of Graph Neural Networks in random graphs classification task. Journal of Complex Networks. 2026;14(2):cnag006. doi: \href{https://doi.org/10.1093/comnet/cnag006}{\nolinkurl{10.1093/comnet/cnag006}}
\bibitem{cortes2021inconsistency} Cortes C, Lawrence ND. Inconsistency in Conference Peer Review: Revisiting the 2014 NeurIPS Experiment. arXiv:2109.09774 [Preprint]. 2021. doi: \href{https://doi.org/10.48550/arxiv.2109.09774}{\nolinkurl{10.48550/arxiv.2109.09774}}
\bibitem{beygelzimer2023consistency} Beygelzimer A, Dauphin YN, Liang P, Vaughan JW. Has the Machine Learning Review Process Become More Arbitrary as the Field Has Grown? The NeurIPS 2021 Consistency Experiment. arXiv:2306.03262 [Preprint]. 2023. doi: \href{https://doi.org/10.48550/arxiv.2306.03262}{\nolinkurl{10.48550/arxiv.2306.03262}}
\bibitem{lu2023calibrating} Lu Y, Kong Y. Calibrating "Cheap Signals" in Peer Review without a Prior. arXiv:2312.07269 [Preprint]. 2023. doi: \href{https://doi.org/10.48550/arxiv.2312.07269}{\nolinkurl{10.48550/arxiv.2312.07269}}
\bibitem{huang2026lottery} Huang F. The Review Lottery: Benchmarking an Observational Estimator of Peer-Review Noise (ICLR 2017-2025). arXiv:2610.06591 [Preprint]. 2026. doi: \href{https://doi.org/10.48550/arxiv.2610.06591}{\nolinkurl{10.48550/arxiv.2610.06591}}
\bibitem{xu2026comparable} Xu B, Dai X, Yang F, Zhang K. Reviewer Scores Are Not Comparable Across Research Areas in ML Peer Review. arXiv:2607.27209 [Preprint]. 2026. doi: \href{https://doi.org/10.48550/arxiv.2607.27209}{\nolinkurl{10.48550/arxiv.2607.27209}}
\bibitem{gonzalez2024iclr} González-Márquez R, Kobak D. Learning representations of learning representations. arXiv:2404.08403 [Preprint]. 2024. doi: \href{https://doi.org/10.48550/arxiv.2404.08403}{\nolinkurl{10.48550/arxiv.2404.08403}}
\bibitem{yang2025qwen3} Yang A, Li A, Yang B, Zhang B, Hui B, Zheng B, et al. Qwen3 Technical Report. arXiv:2505.09388 [Preprint]. 2025. doi: \href{https://doi.org/10.48550/arxiv.2505.09388}{\nolinkurl{10.48550/arxiv.2505.09388}}
\bibitem{kwon2023vllm} Kwon W, Li Z, Zhuang S, Sheng Y, Zheng L, Yu CH, et al. Efficient Memory Management for Large Language Model Serving with PagedAttention. In: Proceedings of the 29th Symposium on Operating Systems Principles; 2023. p. 611-626. doi: \href{https://doi.org/10.1145/3600006.3613165}{\nolinkurl{10.1145/3600006.3613165}}
\bibitem{fey2019pyg} Fey M, Lenssen JE. Fast Graph Representation Learning with PyTorch Geometric. arXiv:1903.02428 [Preprint]. 2019. doi: \href{https://doi.org/10.48550/arxiv.1903.02428}{\nolinkurl{10.48550/arxiv.1903.02428}}
\bibitem{kokhlikyan2020captum} Kokhlikyan N, Miglani V, Martin M, Wang E, Alsallakh B, Reynolds J, et al. Captum: A unified and generic model interpretability library for PyTorch. arXiv:2009.07896 [Preprint]. 2020. doi: \href{https://doi.org/10.48550/arxiv.2009.07896}{\nolinkurl{10.48550/arxiv.2009.07896}}
\bibitem{cohen1960kappa} Cohen J. A Coefficient of Agreement for Nominal Scales. Educational and Psychological Measurement. 1960;20(1):37-46. doi: \href{https://doi.org/10.1177/001316446002000104}{\nolinkurl{10.1177/001316446002000104}}
\bibitem{reimers2019sbert} Reimers N, Gurevych I. Sentence-BERT: Sentence Embeddings using Siamese BERT-Networks. In: Proceedings of the 2019 Conference on Empirical Methods in Natural Language Processing and the 9th International Joint Conference on Natural Language Processing (EMNLP-IJCNLP); 2019. p. 3980-3990. doi: \href{https://doi.org/10.18653/v1/d19-1410}{\nolinkurl{10.18653/v1/d19-1410}}
\bibitem{song2020mpnet} Song K, Tan X, Qin T, Lu J, Liu TY. MPNet: Masked and Permuted Pre-training for Language Understanding. arXiv:2004.09297 [Preprint]. 2020. doi: \href{https://doi.org/10.48550/arxiv.2004.09297}{\nolinkurl{10.48550/arxiv.2004.09297}}
\bibitem{guo2017calibration} Guo C, Pleiss G, Sun Y, Weinberger KQ. On Calibration of Modern Neural Networks. arXiv:1706.04599 [Preprint]. 2017. doi: \href{https://doi.org/10.48550/arxiv.1706.04599}{\nolinkurl{10.48550/arxiv.1706.04599}}
\bibitem{efron1994bootstrap} Efron B, Tibshirani RJ. An Introduction to the Bootstrap. Chapman and Hall/CRC; 1994. doi: \href{https://doi.org/10.1201/9780429246593}{\nolinkurl{10.1201/9780429246593}}
\bibitem{delong1988auc} DeLong ER, DeLong DM, Clarke-Pearson DL. Comparing the Areas under Two or More Correlated Receiver Operating Characteristic Curves: A Nonparametric Approach. Biometrics. 1988;44(3):837. doi: \href{https://doi.org/10.2307/2531595}{\nolinkurl{10.2307/2531595}}
\bibitem{aickin1996holm} Aickin M, Gensler H. Adjusting for multiple testing when reporting research results: the Bonferroni vs Holm methods. American Journal of Public Health. 1996;86(5):726-728. doi: \href{https://doi.org/10.2105/ajph.86.5.726}{\nolinkurl{10.2105/ajph.86.5.726}}
\end{thebibliography}}
